%%%%%%%%%%%%%%%%%%%%%%%%%%%%%%%%%%%%%%%%%%%%%%%%%%%%%%%%%%%%%%%%%%%%%%
% Lớp 12 | Phương trình mặt phẳng, đường thẳng, mặt cầu | Phương trình đường thẳng trong không gian
% Dạng 2H5?2-2: Xác định cặp véc-tơ chỉ phương, véc-tơ pháp tuyến
% Số câu: 44  (sinh tự động bởi tach_cau.py; từ nay thầy sửa trực tiếp file này)
%%%%%%%%%%%%%%%%%%%%%%%%%%%%%%%%%%%%%%%%%%%%%%%%%%%%%%%%%%%%%%%%%%%%%%

%%%%%%%%%%%%%%%% TRẮC NGHIỆM 4 PHƯƠNG ÁN %%%%%%%%%%%%%%%%

% nguồn: Phan1_Lop_12_Hinh_hoc_Chuong_5
\begin{ex}%[2H5N2-2]%[ma:90c45754]
	Trong không gian $Oxyz$, cho đường thẳng $\Delta \colon \dfrac{x-1}{3}=\dfrac{y+2}{-5}=\dfrac{z-4}{1}$. Vectơ nào sau đây là  một vectơ chỉ phương của đường thẳng $\Delta$?
	\choice{$\vec{u}=(1;-2;4)$} 
		 {\True $\vec{u}=(3;-5;1)$}
		 {$\vec{u}=(-1;2;-4)$}
		 {$\vec{u}=(3;-5;4)$}
	 \loigiai{
	Một vectơ chỉ phương của  đường thẳng $\Delta$ là $\vec{u}=(3;-5; 1)$.}
\end{ex}

% nguồn: Phan1_Lop_12_Hinh_hoc_Chuong_5
\begin{ex}%[2H5N2-2]%[ma:89054c73]
	Trong không gian $Oxyz$, cho đường thẳng $d \colon  \dfrac{x-3}{2}=\dfrac{y-4}{-5}=\dfrac{z+1}{3}$. Vectơ nào dưới đây là một vectơ chỉ phương của $d$?
	\choice{$\vec{a}=(3;4;-1)$} 
		 {\True $\vec{a}=(2;-5;3)$}
		{$\vec{a}=(-3;-4;1)$} 
		{$\vec{a}=(2;-5;-3)$} 
	\loigiai{ 
	Một vectơ  chỉ phương của đường thẳng $d \colon \dfrac{x-3}{2}=\dfrac{y-4}{-5}=\dfrac{z+1}{3}$ là $\vec{a}=(2;-5;3)$.}
\end{ex}

% nguồn: Phan1_Lop_12_Hinh_hoc_Chuong_5
\begin{ex}%[2H5N2-2]%[ma:3e4260c0]
	Trong không gian $Oxyz$, cho đường thẳng $d\colon \dfrac{x-3}{2}=\dfrac{y-4}{-5}=\dfrac{z+1}{3}$. Véc-tơ nào dưới đây là một  véc-tơ chỉ phương của $d$?
	\choice 
		 {$\overrightarrow{a}=(3;4;-1)$}
		 {\True $\overrightarrow{a}=(2;-5;3)$}
		{$\overrightarrow{a}=(-3;-4;1)$} 
		 {$\overrightarrow{a}=(2;-5;-3)$}
	\loigiai{ 
		Véc-tơ chỉ phương của $d$ là  $\overrightarrow{a}=(2;-5;3)$.
	} 
\end{ex}

% nguồn: Phan1_Lop_12_Hinh_hoc_Chuong_5
\begin{ex}%[2H5N2-2]%[ma:f16e9345]
	Trong không gian $Oxyz$, cho đường  thẳng $\heva{&x=2+3t\\&y=-1+t\\&z=5-2t}$. Vectơ nào sau đây là một vectơ chỉ phương của đường thẳng $\Delta$?
	\choice{$\vec{u}=(2;-1;5)$} 
		 {$(2;-1;-2)$}
		{$\vec{u}=(3;1;5)$} 
		 {\True $\vec{u}=(3;1;-2)$}
	\loigiai{ 
	Một vectơ chỉ phương của đường thẳng $\Delta$ là $\vec{u}=(3;1;-2)$.} 
\end{ex}

% nguồn: Phan1_Lop_12_Hinh_hoc_Chuong_5
\begin{ex}%[2H5N2-2]%[ma:d30a6e97]
	Trong không gian với hệ tọa  độ $Oxyz$, vectơ nào là một vectơ chỉ phương của đường thẳng có phương trình: $\dfrac{x-1}{3}=\dfrac{3 y}{2}=\dfrac{3-z}{1}$?
	\choice{$\vec{a}=\left(3; \dfrac{3}{2}; 1\right)$} 
		{$\vec{a}=(9; 2;  3)$}
		{$\vec{a}=(3;  2; 1)$}
		{\True $\vec{a}=\left(3;  \dfrac{2}{3};-1\right)$}
	\loigiai{ 
		Ta có:  $\dfrac{x-1}{3}=\dfrac{3 y}{2}=\dfrac{3-z}{1} \Rightarrow \dfrac{x-1}{3}=\dfrac{y}{\dfrac{2}{3}}=\dfrac{z-3}{-1}$.\\
	Một vectơ  chỉ phương của đường thẳng có phương trình là $\vec{a}=\left(3; \dfrac{2}{3};-1\right)$.}
\end{ex}

% nguồn: Phan1_Lop_12_Hinh_hoc_Chuong_5
\begin{ex}%[2H5N2-2]%[ma:b69a8d3d]
	Trong không gian $ Oxyz$, véc-tơ nào sau đây~là một véc-tơ chỉ phương của  trục $Ox$?
	\choice 
		{\True ${\vec  i}=\left(1;0;0\right)$}
		{$\vec  j=\left(0;1;0\right)$}
		 {${\vec k}=\left(0;0;1\right)$}
		 {${\vec u}=\left(1;1;1\right)$}
	\loigiai{ 
		véc-tơ chỉ phương của trục  $Ox$~là ${\vec i}=\left(1;0;0\right)$.
	} 
\end{ex}

% nguồn: Phan1_Lop_12_Hinh_hoc_Chuong_5
\begin{ex}%[2H5N2-2]%[ma:73ef33b6]
	Trong không gian $ Oxyz$, véc-tơ  nào sau đây~là một véc-tơ chỉ phương của trục $Oy$ ?
	\choice 
		{${\vec i}=\left(1;0;0\right)$} 
		 {\True $\vec j=\left(0;1;0\right)$}
		 {${\vec k}=\left(0;0;1\right)$}
		{${\vec u}=\left(1;1;1\right)$} 
	 \loigiai{
		Véc-tơ  chỉ phương của trục $Oy$~là $\vec j=\left(0;1;0\right)$.
	} 
\end{ex}

% nguồn: Phan1_Lop_12_Hinh_hoc_Chuong_5
\begin{ex}%[2H5N2-2]%[ma:74cde8dd]
	Trong không gian $Oxyz$, tọa độ nào sau đây là tọa độ của một  vectơ chỉ phương của đường thẳng
	$\Delta \colon \heva{&x=2+4t\\&y=1-6t\\&z=9t}(t\in\mathbb{R})$? 
	 \choice{$\vec{a}=\left(\dfrac{1}{3};\dfrac{1}{2};\dfrac{3}{4}\right)$}
		{$\vec{a}=(2;1;0)$} 
		{$\vec{a}=(4;-6;0)$} 
		{\True  $\vec{a}=\left(\dfrac{1}{3};-\dfrac{1}{2};\dfrac{3}{4}\right)$}
	\loigiai{ 
	Một vectơ chỉ phương của đường thẳng $\Delta \colon \heva{&x=2+4t\\&y=1-6t\\&z=9t}(t\in\mathbb{R})$ là $\vec{a}=\left(\dfrac{1}{3};-\dfrac{1}{2};\dfrac{3}{4}\right)$.} 
\end{ex}

% nguồn: Phan1_Lop_12_Hinh_hoc_Chuong_5
\begin{ex}%[2H5N2-2]%[ma:ba7d05d3]
	Trong không gian $Oxyz$, tọa độ  một véc-tơ chỉ phương của đường thẳng có phương trình $d\colon\dfrac{x-2}{-3}=\dfrac{y+1}{1}=\dfrac{z}{2}$ là
	\choice 
		 {$(2;-1;0)$}
		{\True  $(-3;1;2)$}
		{$(-2;1;0)$} 
		{$(3;1;-2)$} 
	\loigiai{Véc-tơ  chỉ phương của đường thẳng $d$ là $\vec{u}=(-3;1;2)$.}
\end{ex}

% nguồn: Phan1_Lop_12_Hinh_hoc_Chuong_5
\begin{ex}%[2H5N2-2]%[ma:acfb0d24]
	Trong  không gian $Oxyz$, cho đường thẳng $d$ vuông góc với mặt phẳng $(\alpha)$ có phương trình $x+2z+3=0$. Một vectơ chỉ phương của $d$ là
	\choice{\True  $\vec{a}=(1;0;2)$}
		{$\vec{a}=(2;-1;0)$} 
		{$\vec{a}=(1;2;3)$} 
		{$\vec{a}=(2;0;-1)$} 
	 \loigiai{
		Đường thẳng $d$ vuông góc với mặt phẳng $(\alpha) \colon x+2z+3=0$ có một  vectơ chỉ phương là $\vec{a}=(1;0;2)$.
	 }
\end{ex}

% nguồn: Phan1_Lop_12_Hinh_hoc_Chuong_5
\begin{ex}%[2H5N2-2]%[ma:b5c1e3be]
	Trong  không gian $Oxyz$, cho đường thẳng $d\colon \heva{&x=2+2t\\&y=1+3t\\&z=-1}$. Một VTCP của đường thẳng $d$ là
	 \choice
		 {$(2;3;-1)$}
		{\True  $(2;3;0)$}
		 {$(2;1;-1)$}
		{$(2;1;  0)$}
	 \loigiai{
		Vì $d\colon \heva{&x=2+2t\\&y=1+3t\\&z=-1}$ nên một VTCP của đường thẳng $d$ là $(2;3;0)$. 
	 }
\end{ex}

% nguồn: Phan1_Lop_12_Hinh_hoc_Chuong_5
\begin{ex}%[2H5N2-2]%[ma:c7417ec3]
	Trong không gian $Oxyz$, cho hai điểm  $M(3; 0; 0)$ và $N(0; 0; 1)$. Vectơ nào dưới đây là một vectơ chỉ phương của đường thẳng $MN$?
	 \choice
		 {$\vec{d}=(-1; 1; 2)$}
		{$\vec{c}=(1; 2;  2)$}
		 {$\vec{b}=(3; 0; 1)$}
		{\True  $\vec{a}=(-3; 0; 1)$}
	 \loigiai{
		Ta có $\overrightarrow{MN}=(-3;0;1)$ nên vectơ $\vec{a}=(-3; 0; 1)$ là một vectơ chỉ  phương của $MN$.
	} 
\end{ex}

% nguồn: Phan1_Lop_12_Hinh_hoc_Chuong_5
\begin{ex}%[2H5N2-2]%[ma:32a93946]
	Trong không gian $Oxyz$, đường  thẳng $d \colon \heva{&x=1-t \\&y = 3+ 2t \\&z=t}$ có một vec-tơ chỉ phương là
	 \choice
		 {$\vec{u}_{1} =(-1;2;0)$}
		{$\vec{u}_{2} = (1;3;1)$} 
		{$\vec{u}_{3} =  (1;2;1)$}
		{\True  $\vec{u}_{4} =(-1;2;1)$}
	 \loigiai{
		 Một vectơ chỉ phương của đường thẳng $d$ là $\vec{u}_{4} =(-1;2;1)$
	} 
\end{ex}

% nguồn: Phan1_Lop_12_Hinh_hoc_Chuong_5
\begin{ex}%[2H5N2-2]%[ma:b31a4b0e]
	Trong không gian $Oxyz$, cho đường thẳng $d\colon \dfrac{x-3}{2}=\dfrac{y-4}{-5}=\dfrac{z+1}{3}$. Vectơ nào dưới đây là một vectơ  chỉ phương của $d$?
	\choice 
		{$\vec{u_2}=(2;4;-1)$} 
		 {$\vec{u_4}=(3;4;-1)$}
		{\True$\vec{u_1}=(2;-5;3)$} 
		 {$\vec{u_3}=(2;5;3)$}
	 \loigiai{
		Vectơ  chỉ phương của đường thẳng $d$ là $\vec{u}=(2;-5;3)$.
	} 
\end{ex}

% nguồn: Phan1_Lop_12_Hinh_hoc_Chuong_5
\begin{ex}%[2H5N2-2]%[ma:24b846e1]
	Trong không  gian $Oxyz$, vectơ nào sau đây là vectơ chỉ phương của đường thẳng $\Delta \colon \dfrac{x-5}{8}=\dfrac{y-9}{6}=\dfrac{z-12}{3}$.
	\choice 
		{$\overrightarrow{u}_2=(8;  6;-3)$}
		{$\overrightarrow{u}_4=(5;  9; 12)$}
		{\True $\overrightarrow{u}_1=(8; 6;  3)$}
		{$\overrightarrow{u}_3=(-8;  6;-3)$}
	\loigiai{Vectơ  chỉ phương của đường thẳng $\Delta$ là $\overrightarrow{u}_1=(8; 6; 3)$.}
\end{ex}

% nguồn: Phan1_Lop_12_Hinh_hoc_Chuong_5
\begin{ex}%[2H5N2-2]%[ma:6d3e11dd]
	. 
	Trong không gian $Oxyz$, vectơ chỉ phương của đường thẳng $d : \begin{cases} x=1+2t \\ y = 3t \\ z=1+t
\end{cases} (t\in\mathbb{R})$ là
\choice 
	{\True $\vec{u} =  (2;3;1)$}
	{$\vec{u} = (1;3;1)$} 
	{$\vec{u} = (3;2;1)$} 
	{$\vec{u} = (2;3;  -1)$}
\loigiai{ 
 Vectơ chỉ phương của đường thẳng $d \colon \heva{&x=1+2t \\&y = 3t \\&z=1+t}(t\in\mathbb{R})$ là $\vec{u} = (2;3;1)$.}
\end{ex}

% nguồn: Phan1_Lop_12_Hinh_hoc_Chuong_5
\begin{ex}%[2H5N2-2]%[ma:17fd84ee]
	Trong không gian $Oxyz$, cho đường thẳng $\Delta\colon\heva{& x=1+t\\ &y=-2+2t\\  &z=3-t}$.
	 Véc-tơ nào sau đây~là một véc-tơ chỉ phương của đường thẳng $\Delta$?
	 \choice
		{$\vec{u}_1=(1;-2;3)$} 
		{\True $\vec{u}_2=(1;2;-1)$} 
		{$\vec{u}_3=(1;2;1)$} 
		 {$\vec{u}_4=(2;0;2)$}
	 \loigiai{
		Đường thẳng $\Delta$  có một véc-tơ chỉ phương~là $\vec{u}_2=(1;2;-1)$.
	} 
\end{ex}

% nguồn: Phan1_Lop_12_Hinh_hoc_Chuong_5
\begin{ex}%[2H5N2-2]%[ma:a365e267]
	Trong không gian  $Oxyz$, cho đường thẳng $d\colon \dfrac{x+1}{1} = \dfrac{y - 3}{-1} = \dfrac{z - 2}{2}$. Vectơ nào dưới đây là một vectơ chỉ phương của $d$?
	\choice 
		{$\overrightarrow{u}_1 =  (1;1;2)$}
		{$\overrightarrow{u}_4 = (1;-3;2)$} 
		{$\overrightarrow{u}_3  = (1;-3;-2)$}
		{\True $\overrightarrow{u}_2 =  (1;-1;2)$}
	\loigiai{Đường thẳng $d\colon\dfrac{x + 1}{1} = \dfrac{y - 3}{-1} = \dfrac{z - 2}{2}$ có một vectơ chỉ phương là $\overrightarrow{u} = (1;-1;2)$. 
	} 
\end{ex}

% nguồn: Phan1_Lop_12_Hinh_hoc_Chuong_5
\begin{ex}%[2H5N2-2]%[ma:cec7a373]
	Trong không gian với  hệ tọa độ $Oxyz$, cho đường thẳng $d \colon$
	$\begin{cases}
	x = 1+2t \\ 
	 y = -3 - t \\
	z = 2 -  3t
\end{cases}$ ($t \in \mathbb{R}$).
 Một vectơ chỉ phương của đường thẳng $d$ là
 \choice
	{$\overrightarrow{u} = (-1; 3; -2)$} 
	 {$\overrightarrow{u} = (2; 1; -3)$}
	 {$\overrightarrow{u} = (-2; 1; 3)$}
	 {\True $\overrightarrow{u} = (2; -1; -3)$}
\loigiai{ 
	Từ phương trình tham số, vectơ chỉ phương của $d$  là $\overrightarrow{u} = (2, -1, -3)$.
 }
\end{ex}

% nguồn: Phan1_Lop_12_Hinh_hoc_Chuong_5
\begin{ex}%[2H5N2-2]%[ma:7e36c966]
	Trong không gian $Oxyz$, đường thẳng $d\colon \dfrac{x-1}{2}=\dfrac{y+1}{-1}=\dfrac{z}{2}$ có vectơ chỉ phương  là
	 \choice
		{\True  $(2;-1;2)$}
		 {$(1;-1;1)$}
		 {$(3;-1;0)$}
		 {$(1;-1;0)$}
	\loigiai{ 
		Đường thẳng $d\colon \dfrac{x-1}{2}=\dfrac{y+1}{-1}=\dfrac{z}{2}$  có vectơ chỉ phương $(2;-1;2)$.
	 }
\end{ex}

% nguồn: Phan1_Lop_12_Hinh_hoc_Chuong_5
\begin{ex}%[2H5N2-2]%[ma:368c566a]
	Trong không gian $Oxyz$, vectơ nào  sau đây~là một vectơ chỉ phương của đường thẳng $\Delta\colon \dfrac{x-5}{8}=\dfrac{y-9}{6}=\dfrac{z-12}{3}$?
	\choice 
		{\True $\vec{u}_1=(8;  6; 3)$}
		{$\vec{u}_2=(8; 6;-3)$} 
		{$\vec{u}_3=(-8;  6;-3)$}
		{$\vec{u}_4=(5; 9;  12)$}
	 \loigiai{
		Vectơ chỉ phương của  đường thẳng $\Delta\colon \dfrac{x-5}{8}=\dfrac{y-9}{6}=\dfrac{z-12}{3}$~là $\vec{u}_1=(8; 6; 3)$.
	} 
\end{ex}

% nguồn: Phan1_Lop_12_Hinh_hoc_Chuong_5
\begin{ex}%[2H5N2-2]%[ma:484d8259]
	Đường thẳng $d$ đi qua  $M(2; 0;-1)$ và có véc-tơ chỉ phương $\overrightarrow{a}=(2;-3; 1)$ có phương trình
	 \choice
		{$\heva{&x=4+2t\\&y=-3t\\&z=2+t}$} 
		 {$\heva{&x=-2+2t\\&y=-3t\\&z=1+t}$}
		{\True  $\heva{&x=2+2t\\&y=-3t\\&z=-1+t}$}
		 {$\heva{&x=-2+4t\\&y=-6t\\&z=1+2t}$}
	\loigiai{Đường thẳng $d$ đi qua $M(2; 0;-1)$ và có véc-tơ chỉ phương $\overrightarrow{a}=(2;-3; 1)$ là  $d\colon\heva{&x=2+2t\\&y=-3t\\&z=-1+t.}$
	} 
\end{ex}

% nguồn: Phan1_Lop_12_Hinh_hoc_Chuong_5
\begin{ex}%[2H5N2-2]%[ma:b0eee9e0]
	Trong không gian $Oxyz$, đường thẳng $\heva{&x=2-t\\&y=1+2t\\&z=3+t}$ có một  vectơ chỉ phương là
	 \choice{\True $\vec{a}=(-1;2;1)$}
		{$\vec{a}=(2;1;3)$} 
		 {$\vec{a}=(1;2;3)$}
		{$\vec{a}=(2;0;-1)$} 
	 \loigiai{
		Đường thẳng $\heva{&x=2-t\\&y=1+2t\\&z=3+t}$ có một vectơ chỉ phương  là $\vec{a}=(-1;2;1)$.
	} 
\end{ex}

% nguồn: Phan1_Lop_12_Hinh_hoc_Chuong_5
\begin{ex}%[2H5N2-2]%[ma:061afc34]
	Trong không gian $Oxyz$, cho đường thẳng $d\colon \dfrac{x-1}{-1}=\dfrac{y}{2}=\dfrac{z-5}{3}$. véc-tơ nào dưới đây là một véc-tơ  chỉ phương của $d$?
	 \choice
		{$\overrightarrow{u}_3=(1;2;-3)$} 
		 {$\overrightarrow{u}_1=(1;-2;3)$}
		{\True  $\overrightarrow{u}_4=(1;-2;-3)$}
		 {$\overrightarrow{u}_2=(1;2;3)$}
	\loigiai{ 
		Một véc-tơ chỉ  phương của $d$
	 là $\overrightarrow{u}_4=(1;-2;-3)$.}
\end{ex}

% nguồn: Phan1_Lop_12_Hinh_hoc_Chuong_5
\begin{ex}%[2H5N2-2]%[ma:d01f349f]
	Trong không gian $Oxyz$, cho đường thẳng $d  \colon \dfrac{x-3}{4} = \dfrac{y+1}{-2} = \dfrac{z+2}{3}$. Vectơ nào dưới đây là một vectơ chỉ phương của $d$?
	 \choice
		{$\overrightarrow{u}_3 =  (3;-1;-2)$}
		{$\overrightarrow{u}_4 =  (4;2;3)$}
		{\True $\overrightarrow{u}_2 = (4;-2;3)$} 
		{$\overrightarrow{u}_1 =  (3;1;2)$}
	\loigiai{ 
		Một vectơ chỉ  phương của đường thẳng $d$ là $\overrightarrow{u}_2 = (4;-2;3)$.
	} 
\end{ex}

% nguồn: Phan1_Lop_12_Hinh_hoc_Chuong_5
\begin{ex}%[2H5N2-2]%[ma:f2b39cff]
	Trong không gian $Oxyz$, đường thẳng đi qua hai điểm  $A(-1;1;0)$ và $B(3;2;-1)$ có một véc-tơ chỉ phương~là
	\choice 
		{$\vec{a}=(2;1;1)$} 
		{\True  $\vec{u}=(4;1;-1)$}
		{$\vec{b}=(2;1;-1)$} 
		 {$\vec{v}=(4;1;1)$}
	 \loigiai{
		Đường thẳng  đi qua hai điểm $A(-1;1;0)$ và $B(3;2;-1)$ có một véc-tơ chỉ phương~là $\overrightarrow{AB}=(4;1;-1)$.
	 }
\end{ex}

% nguồn: Phan1_Lop_12_Hinh_hoc_Chuong_5
\begin{ex}%[2H5N2-2]%[ma:0eda9819]
	Trong không gian $ Oxyz$, véc-tơ nào sau đây~là một véc-tơ  chỉ phương của trục $Oz$?
	 \choice
		 {${\vec i}=\left(1;0;0\right)$}
		{${\vec  j}=\left(0;1;0\right)$}
		{\True $\vec k=\left(0;0;1\right)$} 
		 {${\vec u}=\left(1;1;1\right)$}
	 \loigiai{
		véc-tơ chỉ phương  của trục $Oz$~là $\vec k=\left(0;0;1\right)$.
	} 
\end{ex}

% nguồn: Phan1_Lop_12_Hinh_hoc_Chuong_5
\begin{ex}%[2H5N2-2]%[ma:4061e6ba]
	Trong không gian  $Oxyz$, cho đường thẳng $\Delta:\dfrac{x-1}{2}=\dfrac{y-5}{1}=\dfrac{z+2}{-5}$. Vectơ nào sau đây là vectơ chỉ phương của đường thẳng $\Delta$?
	 \choice
		{\True  $\vec{u}=(2;1;-5)$}
		 {$\vec{u}=(1;5;-2)$}
		 {$\vec{u}=(-1;-5;2)$}
		 {$\vec{u}=(2;-1;5)$}
	\loigiai{ 
		$\vec{u}=(2;1;-5)$. 
	 }
\end{ex}

% nguồn: Phan1_Lop_12_Hinh_hoc_Chuong_5
\begin{ex}%[2H5N2-2]%[ma:e6d1d0b2]
	Trong không gian $Oxyz$, cho đường thẳng $(d)\colon\dfrac{x-1}{2}=\dfrac{y+2}{-1}=\dfrac{z-3}{1}$. Vectơ nào sau  đây là một vectơ chỉ phương của đường thẳng $(d)$ ?
	\choice 
		 {$\overrightarrow{u}_1=(-1;2;-3)$}
		 {$\overrightarrow{u}_2=(1;-2;3)$}
		 {$\overrightarrow{u}_3=(2;1;1)$}
		{\True  $\overrightarrow{u}_4=(2;-1;1)$}
	\loigiai{ 
		 Vectơ chỉ phương của đường thẳng $(d)$ là $\overrightarrow{u}_4=(2;-1;1)$.
	 }
\end{ex}

% nguồn: Phan1_Lop_12_Hinh_hoc_Chuong_5
\begin{ex}%[2H5N2-2]%[ma:e07a77b5]
	Trong không gian với hệ tọa độ $Oxyz$, cho đường thẳng $d\colon \heva{&x=0\\&y=t\\&z=2-t}$. Vectơ nào dưới đây là vecto  chỉ phương của đường thẳng $d$ ?
	 \choice
		{\True $\overrightarrow{u}=(1; 0;-1)$} 
		{$\overrightarrow{u}=(0;  1;-1)$}
		{$\overrightarrow{u}=(0; 0; 2)$} 
		{$\overrightarrow{u}=(0; 1;  2)$}
	\loigiai{Vectơ chỉ  phương của đường thẳng $d$ là $\overrightarrow{u}=(0; 1;-1)$.}
\end{ex}

% nguồn: Phan1_Lop_12_Hinh_hoc_Chuong_5
\begin{ex}%[2H5H2-2]%[ma:c652a39d]
	Trong không gian với hệ tọa độ $Oxyz$, cho hai mặt phẳng $(P)\colon 2x+y-z-1=0$ và $(Q)\colon x-2y+z-5=0$. Khi đó giao tuyến của hai mặt phẳng $(P)$ và $(Q)$ có  một vectơ chỉ phương là
	\choice 
		 {$\vec{u}_1=(1;-2;1)$}
		{$\vec{u}_2=(2;1;-1)$} 
		 {\True $\vec{u}_3=(1;3;5)$}
		 {$\vec{u}_4=(-1;3;-5)$}
	\loigiai{ 
		Gọi $d=(P)\cap(Q)$. Khi đó một vectơ chỉ phương của $d$ là  $$\vec{u}=\left[\vec{n}_{P},\vec{n}_{Q}\right]=(-1;-3;-5).$$
		Vậy $\vec{u}_3=(1;3;5)$ là một vectơ chỉ phương  của $d$.
	 }
\end{ex}

% nguồn: Phan1_Lop_12_Hinh_hoc_Chuong_5
\begin{ex}%[2H5H2-2]%[ma:53de4e98]
	Trong không gian $(Oxyz)$,  cho đường thẳng $d\colon \dfrac{x-1}{2}=\dfrac{y-2}{-3}=\dfrac{z+1}{2}$. Véc-tơ nào sau đây là véc-tơ chỉ phương của đường thẳng $d$?
	 \choice
		 {$\overrightarrow{u}=(2;3;2)$}
		{$\overrightarrow{u}=(1;1;3)$} 
		{$\overrightarrow{u}=(-2;-3;-2)$} 
		{\True $\overrightarrow{u}=(2;-3;2)$} 
	 \loigiai{
		Vì $d\colon \dfrac{x-1}{2}=\dfrac{y-2}{-3}=\dfrac{z+1}{2}$ nên véc-tơ chỉ phương của đường thẳng $d$ là $\overrightarrow{u}=(2;-3;2)$. 
	 }
\end{ex}

% nguồn: Phan1_Lop_12_Hinh_hoc_Chuong_5
\begin{ex}%[2H5H2-2]%[ma:c255ac6e]
	Trong không gian $ Oxyz$, cho hình bình hành $ABCD$. Véc-tơ  nào sau đây~là một véc-tơ chỉ phương của đường thẳng $AB$ ?
	 \choice
		{$\overrightarrow{{BD}}$} 
		 {$\overrightarrow{{AC}}$}
		{$\overrightarrow{{AD}}$} 
		{\True  $\overrightarrow{{CD}}$}
	\loigiai{ 
		Véc-tơ chỉ phương  của đường thẳng $AB$~là $\overrightarrow{{CD}}$.
	 }
\end{ex}

% nguồn: Phan1_Lop_12_Hinh_hoc_Chuong_5
\begin{ex}%[2H5H2-2]%[ma:0e5ddd6b]
	Trong không  gian với hệ trục tọa độ $Oxyz$, cho đường thẳng $d\colon \heva{&x=-1+t\\&y=1+3t\\&z=-3-t}$. Điểm nào dưới đây thuộc $d$?
	 \choice
		{$M(1;-1; 3)$} 
		 {$N(0;-2;-2)$}
		{\True $Q(0;  4;-4)$}
		{$P(0;-4;  4)$}
	\loigiai{ 
		Điểm  $Q(0; 4;-4)$ thuộc $d$.
	 }
\end{ex}

% nguồn: Phan1_Lop_12_Hinh_hoc_Chuong_5
\begin{ex}%[2H5H2-2]%[ma:c49e4309]
	Cho đường thẳng $d\colon \dfrac{x-2}{-1}=\dfrac{y-1}{2}=\dfrac{z+3}{1}$. Vectơ nào dưới đây là vectơ chỉ phương của đường thẳng $d$ 
	 \choice
		 {$\overrightarrow{u}_1=(2; 1;-3)$}
		{$\overrightarrow{u}_2=(-2;-1; 3)$} 
		{\True $\overrightarrow{u}_3=(-1; 2;  1)$}
		{$\overrightarrow{u}_4=(-1;  2;-1)$}
	\loigiai{ 
		$d$ có vectơ chỉ phương là  $\overrightarrow{u}_3=(-1; 2; 1)$.
	} 
\end{ex}

% nguồn: Phan1_Lop_12_Hinh_hoc_Chuong_5
\begin{ex}%[2H5H2-2]%[ma:bf00e67f]
	Trong không  gian với hệ tọa độ $Oxyz$, cho $\overrightarrow{OA}=2\vec{i}+3\vec{j}-5\vec{k}$; $\overrightarrow{OB}=-2\vec{j}-4\vec{k}$. Một vectơ chỉ phương của đường thẳng $AB$ là
	\choice{$\vec{a}=(-2;-5;-1)$} 
		{\True $\vec{a}=(2;5;-1)$} 
		 {$\vec{a}=(2;5;-9)$}
		{$\vec{a}=(2;1;-9)$} 
	\loigiai{ 
	Một vectơ  chỉ phương của đường thẳng $AB$ là $\vec{a}=(2;5;-1)$.}
\end{ex}

% nguồn: Phan1_Lop_12_Hinh_hoc_Chuong_5
\begin{ex}%[2H5H2-2]%[ma:0e0924b1]
	Trong  không gian với hệ tọa độ $Oxyz$, cho hai mặt phẳng $(P)\colon 3x-2y+2z-5=0$ và $(Q)\colon 4x+5y-z+1=0$. Biết hai điểm $A$, $B$ phân biệt cùng thuộc giao tuyến của hai mặt phẳng $(P)$ và $(Q)$. Khi đó $\vec{AB}$ cùng phương với vectơ nào sau đây?
	\choice 
		{\True  $\vec{u}_1=(8;-11;-23)$}
		{$\vec{u}_2=(3;-2;2)$} 
		{$\vec{u}_3=(-8;11;-23)$} 
		{$\vec{u}_4=(4;5;-1)$} 
	 \loigiai{
		Ta có $(P)$ có vectơ pháp tuyến là $\vec{n}_{(P)}=(3;-2;2)$ và $(Q)$ có vectơ  pháp tuyến là $\vec{n}_{(Q)}=(4;5;-1)$.\\
		Do $\heva{&AB\subset(P)\\ &AB\subset(Q)}\Rightarrow\heva{&AB\perp\vec{n}_{(P)}\\ &AB\perp\vec{n}_{(Q)}}$ nên  đường thẳng $AB$ có một vectơ chỉ phương là
		$$\vec{u}=\left[\vec{n}_{(P)},\vec{n}_{(Q)}\right]=(-8;11;23).$$ 
		Do $\vec{AB}$ cũng là một vectơ chỉ phương của $AB$ nên  $\vec{AB}$ cùng phương với vectơ $\vec{u}_1=(8;-11;-23)$.
	} 
\end{ex}


%%%%%%%%%%%%%%%% ĐÚNG/SAI %%%%%%%%%%%%%%%%

% nguồn: Phan1_Lop_12_Hinh_hoc_Chuong_5
\begin{ex}%[2H5H2-2]%[ma:d6f3f0ec]
	Trong không gian với hệ tọa độ $Oxyz$, cho hai đường thẳng\\ $\Delta_1 \colon \dfrac{x-1}{2}=\dfrac{y-2}{1}=\dfrac{z-3}{-2}$ và $\Delta_2 \colon \dfrac{x-4}{-1}=\dfrac{y-5}{-2}=\dfrac{z-6}{2}$. 
	\choiceTF 
		{Vectơ có tọa độ $(1;2;3)$ là một vectơ chỉ phương của  $\Delta_1$}
		{Vectơ  có tọa độ $(4;5;6)$ là một vectơ chỉ phương của $\Delta_2$}
		{\True  Côsin của góc giữa hai vectơ $\overrightarrow{u_1}=(2;1;-2)$ và $\overrightarrow{u_2}=(-1;-2;2)$ bằng $-\dfrac{8}{9}$}
		{Góc giữa hai đường thẳng $\Delta_1$  và $\Delta_2$ (làm tròn kết quả đến hàng đơn vị của độ) bằng $132^\circ$}
	\loigiai{ 
		\begin{itemchoice}
			 \itemch \textbf{Sai}. \\
			Vì vectơ có tọa độ $(1;2;3)$ không cùng phương với vectơ chỉ phương có tọa độ $(2;1-2)$ của $\Delta_1$. 
			 \itemch \textbf{Sai}. \\
			Vì vectơ có tọa độ $(4;5;6)$  không cùng phương với vectơ chỉ phương có tọa độ $(-1;-2;2)$ của $\Delta_2$.
			 \itemch \textbf{Đúng}. \\
			Ta  có
			$$\cos \left(\overrightarrow{u_1},\overrightarrow{u_2}\right)=\dfrac{\overrightarrow{u_1} \cdot \overrightarrow{u_2}}{\left| \overrightarrow{u_1}  \right| \cdot \left| \overrightarrow{u_2} \right|}=\dfrac{2 \cdot (-1)+1 \cdot (-2)+(-2) \cdot 2}{\sqrt{2^2+1^2+(-2)^2} \cdot \sqrt{(-1)^2+(-2)^2+2^2}}=-\dfrac{8}{9}.$$
			 \itemch Sai. Ta có
			 $$\cos \left(\Delta_1,\Delta_2\right)=\left|\cos \left(\overrightarrow{u_1},\overrightarrow{u_2}\right) \right|=\dfrac{8}{9}.$$
			Do đó góc giữa hai đường thẳng $\Delta_1$ và $\Delta_2$ (làm tròn kết quả đến hàng đơn vị của  độ) bằng $27^\circ$.
		\end{itemchoice}
	 }
\end{ex}


%%%%%%%%%%%%%%%% TRẢ LỜI NGẮN %%%%%%%%%%%%%%%%

% nguồn: Phan2_Lop_12_Hinh_hoc_Chuong_5
\begin{bt}%[2H5H2-2]%[ma:8dd85884]
	Trong không gian với hệ tọa độ $Oxyz$, cho đường thẳng $\Delta \colon \dfrac{x-1}{2}=\dfrac{y-3}{-1}=\dfrac{z+5}{-5}$. Biết rằng $\overrightarrow{u}=(a;b;5)$ là một vectơ chỉ phương của $\Delta$. Giá trị của $a+b$  bằng bao nhiêu?
	 \shortans{-1}
	\loigiai{Đường thẳng $\Delta \colon \dfrac{x-1}{2}=\dfrac{y-3}{-1}=\dfrac{z+5}{-5}$ có một vectơ chỉ phương là $\overrightarrow{v}=(2;-1;-5)$.\\ 
		Mà $\overrightarrow{u}=(a;b;5)$ là một vectơ chỉ  phương của $\Delta$ nên $\overrightarrow{u}$ và $\overrightarrow{v}$ là hai vectơ cùng phương.\\
		Do đó $\dfrac{a}{2}=\dfrac{b}{-1}=\dfrac{5}{-5}=-1 \Rightarrow \heva{&\dfrac{a}{2}=-1\\&\dfrac{b}{-1}=-1} \Rightarrow \heva{&a=-2\\&b=1.}$\\ 
		Do đó  $a=-2$, $b=1$ nên $a+b=-2+1=-1$.
	} 
\end{bt}

% nguồn: Phan2_Lop_12_Hinh_hoc_Chuong_5
\begin{bt}%[2H5H2-2]%[ma:10a4a1c8]
	Trong không gian $Oxyz$, cho hai điểm $A(-1; 2; 4)$ và $B(5; 2;-2)$. Biết $\vec{u}=(a; b; 1)$ là vectơ chỉ phương của đường thẳng $\Delta$ là đường đối xứng với $A B$ qua mặt phẳng $(O y z)$. Tính  $a+b$.
	\shortans{1} 
	\loigiai{ 
		Do $A'(1;2;4)$ và $B'(-5;2;-2)$ là hai điểm đối xứng với $A,B$ qua mặt phẳng $(Oyz)$. \\ 
		Vậy nên đường thẳng $\Delta$ đi qua $A'$ và  $B'$ nên $\Delta$ nhận $\overline{A' B'}=(-6;0;-6)$ là vectơ chỉ phương. \\
	Khi đó  $\Delta$ nhận $\vec{u}=(1;0;1)$ là vecto chỉ phương $\Rightarrow a=1;b=0\Rightarrow a+b=1$.}
\end{bt}

% nguồn: Phan2_Lop_12_Hinh_hoc_Chuong_5
\begin{bt}%[2H5V2-2]%[ma:61ed46d0]
	Trong không gian với hệ trục tọa độ $Oxyz$, cho hai đường thẳng  $\Delta \colon \dfrac{x-1}{2}=\dfrac{y+4}{3}=\dfrac{z-5}{1}$ và $\Delta': \dfrac{x+3}{2}=\dfrac{y+2}{3}=\dfrac{z-3}{1}$ đối xứng với nhau qua mặt phẳng $(P)$. Biết $(P)$ có vectơ pháp tuyến là $\vec{n}=(a;b;3)$. Tính $a+b$.
	\shortans{1} 
	 \loigiai{
		Đường thẳng $\Delta$  qua $A(1;-4;5)$ và nhận $\vec{u}=(2;3;1)$ là vecto chỉ phương.\\
		Đường thẳng $\Delta'$ qua $B(-3;-2;3)$ và nhận $\vec{u}=(2;3;1)$  là vecto chỉ phương $\Rightarrow \Delta \parallel \Delta'$.\\
		Gọi $A'$ là điểm  trên $\Delta'$ sao cho $A A' \perp \Delta$.\\
		Do $A' \in \Delta' \Rightarrow A'(-3+2 t;-2+3 t; 3+t) \Rightarrow \overrightarrow{A A'}=(-4+2t;2+3t;-2+t)$.\\ 
		Do  $AA'\perp\Delta\Rightarrow\overrightarrow{A A'}\cdot\vec{u}=0\Leftrightarrow-8+4t+6+9 t-2+t=0 \Leftrightarrow t=\dfrac{2}{7}$.\\
		$\Rightarrow  A'\left(\dfrac{-17}{7}; \dfrac{-8}{7}; \dfrac{23}{7}\right) \Rightarrow \overrightarrow{A A'}=\left(\dfrac{-24}{7};\dfrac{20}{7};\dfrac{-12}{7}\right)$.\\
	Khi đó $(P)$ nhận $\vec{n}=-\dfrac{7}{4}\vec{AA'}=(6;-5;3)$ là vectơ  pháp tuyến $\Rightarrow a=6$; $b=-5 \Rightarrow a+b=1$.}
\end{bt}


%%%%%%%%%%%%%%%% TỰ LUẬN %%%%%%%%%%%%%%%%

% nguồn: Phan2_Lop_12_Hinh_hoc_Chuong_5
\begin{bt}%[2H5N2-2]%[ma:b1dfa7b3]
	Trong không gian $Oxyz$, xác định một vectơ chỉ phương của đường thẳng dưới  đây:
	\begin{listEX}
		 \item $d \colon \heva{&x=2-t\\&y=1+2t\\&z=3+t}$.
		\item  $d \colon \dfrac{x-1}{2}=\dfrac{y-2}{1}=\dfrac{z+1}{2}$.
		\item $d  \colon \dfrac{x-1}{2}=\dfrac{y-1}{2}=\dfrac{2-z}{1}$.
	\end{listEX}
	\loigiai{ 
		\begin{listEX}
			\item Đường thẳng $d$ có một vectơ chỉ phương  là $\vec{u}=(-1; 2; 1)$.
			\item Đường thẳng $d$ có một  vectơ chỉ phương là $\vec{v}=(2; 1; 2)$.
			\item Ta viết lại $d \colon \dfrac{x-1}{2}=\dfrac{y-1}{2}=\dfrac{2-z}{1} \Leftrightarrow d \colon  \dfrac{x-1}{2}=\dfrac{y-1}{2}=\dfrac{z-2}{-1} \Rightarrow \vec{u_{1}}=(2; 2;-1)$.
		\end{listEX}
	} 
\end{bt}

% nguồn: Phan2_Lop_12_Hinh_hoc_Chuong_5
\begin{bt}%[2H5N2-2]%[ma:dec0ecfd]
	Cho hình hộp $ABCD.A'B'C'D'$. Trong các vectơ có điểm đầu và điểm cuối là các đỉnh của hình hộp, những vectơ nào là  vectơ chỉ phương của đường thẳng $AC$?
	\loigiai{Vì $AC$ song song với $A'C'$ nên  các vectơ chỉ phương của đường thẳng $AC$ thỏa bài toán là các vectơ sau đây: $\overrightarrow{AC}$, $\overrightarrow{CA}$, $\overrightarrow{A'C'}$, $\overrightarrow{C'A'}$.}
\end{bt}

% nguồn: Phan2_Lop_12_Hinh_hoc_Chuong_5
\begin{bt}%[2H5V2-2]%[ma:6705f9bf]
	Trong không gian với hệ toạ độ $Oxyz$, cho đường thẳng $d_1 \colon \heva{&x=1+2t \\ &y=2+t \\ &z=-2+5t}$ và $d_2 \colon \heva{&x=1+s \\ &y=2+s \\ &z=1+3s}$. Viết phương trình đường vuông góc chung của  hai đường thẳng $d_1$, $d_2$.
	 \loigiai{
		Gọi $\Delta$  là đường vuông góc chung của hai đường thẳng $d_1$, $d_2$. \\
		Đường thẳng $d_1$ có vectơ chỉ phương là $\overrightarrow{u}_1=(2;1;5)$, $d_2$  có vectơ chỉ phương là $\overrightarrow{u}_2=(1;1;3)$. \\
		Gọi $A=\Delta  \cap d_1 \Rightarrow A(1+2t;2+t;-2+5t)$. \\
		Gọi $B=\Delta \cap d_2 \Rightarrow B(1+s;2+s;1+3s)$.  \\
		$\overrightarrow{AB}=(s-2t;s-t;3+3s-5t)$.  \\
		Ta có 
		\begin{eqnarray*}
			& &\heva{\overrightarrow{AB} \cdot \overrightarrow{u}_1&=0  \\ \overrightarrow{AB} \cdot \overrightarrow{u}_2&=0} \\
			&\Leftrightarrow &  \heva{2(s-2t)+s-t+5(3+3s-5t)&=0 \\ s-2t+s-t+3(3+3s-5t)&=0} \\
			&\Leftrightarrow & \heva{18s-30t&=-15 \\  11s-18t&=-9} \\
			&\Leftrightarrow  & \heva{s&=0 \\ t&=\dfrac{1}{2}.}
		\end{eqnarray*}
		Suy ra  $B(1;2;1)$ và $\overrightarrow{AB}=\left(-1;-\dfrac{1}{2};\dfrac{1}{2}\right)$. \\
		Chọn $\overrightarrow{u}=2 \cdot \overrightarrow{AB}=(-2;-1;1)$  là một vectơ chỉ phương của $\Delta$. \\
		Phương trình đường thẳng $\Delta$ đi qua $B(1;2;1)$ và  có vectơ chỉ phương $\overrightarrow{u}=(-2;-1;1)$ là
		$\heva{&x=1-2t \\ &y=2-t  \\ &z=1+t}$.
	} 
\end{bt}
